\documentclass[11pt]{article}
\usepackage[margin=1in]{geometry}
\usepackage{booktabs}
\title{DME: SNc vs.\ VTA, within Intact Cells --- Interpretation}
\author{GSE233866, 6-OHDA lesion cohort}
\date{}

\begin{document}
\maketitle

\subsection*{The gap this fills}

\texttt{../snc\_vs\_vta\_combined\_dme\_lesioned/} tests SNc against VTA \emph{during}
lesioning and finds CACNA1D's module dramatically more active in SNc
($p_{\mathrm{adj}}=4.5\times10^{-123}$). That result was originally framed as evidence of
something specifically active in the vulnerable region \emph{during neurodegeneration}. This
network tests the missing counterfactual: is SNc already this different from VTA in the same
animals' \textbf{intact} hemisphere, before any lesioning?

\subsection*{Result: essentially the same, not disease-induced}

CACNA1D's module (turquoise, 1{,}629 genes, kME=0.327) is higher in SNc than VTA with
$\mathrm{avg\_log2FC}=+13.21$, $p_{\mathrm{adj}}\approx0$ (small enough to underflow double
precision) --- both the direction and magnitude essentially match the lesioned comparison
(+13.34, $p_{\mathrm{adj}}=4.5\times10^{-123}$). \textbf{This changes the framing of study \#7
(and \#6): SNc's molecular distinctiveness from VTA for this program is not something that
emerges or intensifies during neurodegeneration --- it is already present, at comparable
strength, in intact tissue.} The honest claim is not "this activates during disease," it's
"this regional difference exists at baseline in this cohort and persists through lesioning."

\subsection*{Does this weaken the therapeutic argument?}

No, but it changes what it's an argument for. A constitutive, baseline SNc-vs-VTA molecular
distinction is arguably a \emph{better} basis for a region-selective drug target than a
disease-induced one: it means the anatomical/molecular substrate for selectivity doesn't
depend on catching the window during active degeneration, and it's consistent with SNc's
selective vulnerability being a pre-existing property of the region rather than purely a
consequence of the disease process. What this result does rule out is describing \#6/\#7's
SNc-vs-VTA gap as itself evidence of a lesioning-specific mechanism --- that gap was already
there. The genuinely lesioning-specific results in this project remain \#6
(lesioned-vs-intact \emph{within} SNc, i.e.\ does SNc's own network change with lesioning ---
yes) and \#8 (the same test within VTA --- yes, but far more weakly). This file establishes
the SNc-vs-VTA axis is a standing baseline feature, not a disease response.

\subsection*{Consistency with the rest of the project}

Turquoise here shares 74\% of its genes with the independently-built
\texttt{../snc\_intact\_network/} blue module, 52\% with
\texttt{../../healthy\_baseline/snc\_network/}'s turquoise, and 51\% with this project's other
pooled SNc-vs-VTA network (\texttt{../snc\_vs\_vta\_combined\_dme\_lesioned/}). This is now the
fifth independently-clustered network in this project recovering substantially the same SNc
gene neighborhood --- the single most consistent, reproducible signal across the whole mouse
side of this work.

\subsection*{Caveats}

As with every DME table here, treat log2FC magnitude as directional only; blue's
$\mathrm{avg\_log2FC}=+24.38$ is a module-eigengene-near-zero artifact (its own
$p_{\mathrm{adj}}=1.0$ confirms it isn't even significant, so this doesn't affect any real
conclusion). $p_{\mathrm{adj}}$ values reported as exactly 0 by R for turquoise/black/yellow/green
have genuinely underflowed double-precision floating point, not been rounded --- treat them as
``smaller than $10^{-300}$-ish,'' not literally zero.

\end{document}
